\documentclass[conference]{IEEEtran}
\usepackage{graphicx}
\usepackage{amsmath}
\usepackage{amssymb}
\usepackage{booktabs}
\usepackage[caption=false,font=footnotesize]{subfig}
\usepackage{cite}
\usepackage{hyperref}
\usepackage{url}

\graphicspath{{./figures/}}

\title{Learning Semantic Features for Image Understanding}

\author{First Author, Second Author, and Third Author}

\begin{document}

\maketitle

\begin{abstract}
We present a novel approach for learning semantic features from images
using deep neural networks. Our method combines convolutional feature extraction
with attention mechanisms to capture both local and global context. Experiments
on standard benchmarks demonstrate significant improvements over baseline methods.
\end{abstract}

\section{Introduction}

Understanding images at a semantic level is a fundamental challenge in computer
vision. Recent advances in deep learning have enabled remarkable progress in
various vision tasks~\cite{alexnet,vggnet,resnet}.

\section{Related Work}

Convolutional neural networks have been the dominant approach for image
understanding since AlexNet~\cite{alexnet}. VGGNet~\cite{vggnet} demonstrated
that deeper networks with smaller filters achieve better performance.
ResNet~\cite{resnet} introduced skip connections to enable training of
very deep networks.

\section{Proposed Method}

Our method consists of three main components: feature extraction, attention
module, and prediction head.

% ---- Subfigure layout (CVPR style) ----
\begin{figure}[t]
  \centering
  \subfloat[Input image]{\includegraphics[width=0.3\textwidth]{fig-input.png}\label{fig:input}}
  \hfill
  \subfloat[Feature map visualization]{\includegraphics[width=0.3\textwidth]{fig-features.png}\label{fig:features}}
  \hfill
  \subfloat[Prediction result]{\includegraphics[width=0.3\textwidth]{fig-prediction.png}\label{fig:pred}}
  \caption{Qualitative results on the validation set. (a) Input image.
    (b) Intermediate feature maps. (c) Final prediction overlay.}
  \label{fig:qual}
\end{figure}

\section{Experiments}

\subsection{Datasets}

We evaluate on COCO-Stuff~\cite{cocostuff} and ADE20K~\cite{ade20k} datasets.

\subsection{Implementation Details}

We use ResNet-50 as the backbone and train with SGD optimizer.
The learning rate is set to 0.01 with polynomial decay.

% ---- Cross-page super-long equation ----
\begin{equation}
\mathcal{J}_{\text{CVPR}} =
\underbrace{\frac{1}{|\mathcal{V}|}\sum_{v\in\mathcal{V}}
  \mathcal{L}_{\text{mse}}\bigl(\hat{I}_v, I_v^\star\bigr)}_{\text{reconstruction}}
+ \lambda_1\underbrace{\sum_{(u,v)\in\mathcal{E}}
  \bigl\|\nabla\phi(\hat{I}_u)-\nabla\phi(\hat{I}_v)\bigr\|^2}_{\text{smoothness}}
+ \lambda_2\underbrace{\sum_{v\in\mathcal{V}}
  \mathcal{R}\bigl(\hat{I}_v\bigr)}_{\text{regularization}} \label{eq:cvpr-energy}
\end{equation}

% ---- Horizontal wide table ----
\begin{table}[t]
\centering
\caption{Quantitative results on COCO-Stuff validation set.}
\label{tab:quant}
\resizebox{\linewidth}{!}{
\begin{tabular}{l|rrrr|rrrr}
  \toprule
  \multirow{2}{*}{Method} & \multicolumn{4}{c|}{mIoU$\uparrow$} &
                            \multicolumn{4}{c}{FPS$\uparrow$} \\
  \cmidrule{2-5} \cmidrule{6-9}
  & S & M & L & Avg & S & M & L & Avg \\
  \midrule
  DeepLabV3+~\cite{deeplabv3} & 31.2 & 41.6 & 54.8 & 42.5 & 8.1 & 8.1 & 8.1 & 8.1 \\
  PSPNet~\cite{pspnet} & 32.8 & 42.6 & 55.2 & 43.6 & 7.2 & 7.2 & 7.2 & 7.2 \\
  UperNet~\cite{upernet} & 33.1 & 43.2 & 56.1 & 44.1 & 9.5 & 9.5 & 9.5 & 9.5 \\
  \textbf{Ours} & \textbf{35.7} & \textbf{46.2} & \textbf{59.1} & \textbf{47.0} &
                 \textbf{15.4} & \textbf{15.4} & \textbf{15.4} & \textbf{15.4} \\
  \bottomrule
\end{tabular}
}
\end{table}

\section{Conclusion}

We presented a novel approach for semantic feature learning that achieves
state-of-the-art results on multiple benchmarks~\cite{darknet,yolov4,solo,condinst}.

\bibliographystyle{IEEEtran}
\bibliography{refs}

\end{document}
